% Generated draft: requires mathematical and source review.
\documentclass[11pt]{article}
\usepackage[margin=25mm]{geometry}
\usepackage{fontspec}
\setmainfont{Latin Modern Roman}
\usepackage{amsmath,amssymb,mathtools,unicode-math}
\setmathfont{Latin Modern Math}
\usepackage{enumitem,needspace,xurl,hyperref}
\setlength{\parindent}{0pt}
\setlength{\parskip}{5pt}
\setlength{\emergencystretch}{3em}
\newcommand{\N}{\mathbb{N}}
\newcommand{\Z}{\mathbb{Z}}
\newcommand{\Q}{\mathbb{Q}}
\newcommand{\R}{\mathbb{R}}
\newcommand{\C}{\mathbb{C}}
\newcommand{\F}{\mathbb{F}}
\newcommand{\A}{\mathbb{A}}
\newcommand{\PP}{\mathbb{P}}
\newcommand{\E}{\mathbb{E}}
\newcommand{\eps}{\varepsilon}
\newcommand{\dd}{\,\mathrm{d}}
\begin{document}
\section*{UnsolvedMath ID: 30002138}
\textbf{Source problem code:} catalogue entry
\subsection{Definitions and mathematical statement}
Let \((X^{n+1},g)\) be a closed Riemannian manifold, meaning compact and without boundary, with strictly negative sectional curvature. Let \(\pi:(\widetilde X,\widetilde g)\to(X,g)\) be its universal Riemannian covering, with \(\widetilde g=\pi^*g\), and write \(d_{\widetilde g}\) for its distance. The use of a universal cover is understood componentwise if connectedness is not implicit.

For a unit-speed geodesic ray \(\gamma:[0,\infty)\to\widetilde X\), define its Busemann function by
\[
 b_\gamma(x)=\lim_{s\to\infty}\bigl(d_{\widetilde g}(x,\gamma(s))-s\bigr).
\]
A horosphere is a level set \(H_{\gamma,t}=b_\gamma^{-1}(t)\), where \(t\in\R\). Equip each horosphere \(H\) with the restricted Riemannian metric \(\widetilde g|_{TH}\) and its intrinsic length distance
\[
 d_H(p,q)=\inf_{\alpha}\operatorname{Length}_{\widetilde g}(\alpha),
\]
where the infimum is over piecewise smooth curves \(\alpha\) in \(H\) joining \(p\) to \(q\). This interprets the source notation \((H,\widetilde g|_H)\) as the induced Riemannian metric.

A map \(f:(H,d_H)\to(K,d_K)\) is a \((\lambda,C)\)-quasi-isometry if \(\lambda\ge1\), \(C\ge0\),
\[
 \lambda^{-1}d_H(p,q)-C\le d_K(f(p),f(q))\le\lambda d_H(p,q)+C
 \qquad(p,q\in H),
\]
and every point of \(K\) has distance at most \(C\) from \(f(H)\).

\textbf{Question.} For every closed negatively curved \((X^{n+1},g)\), are all horospheres in \((\widetilde X,\widetilde g)\) pairwise quasi-isometric? Explicitly, does
\[
 \forall H,K\quad\exists\lambda_{H,K}\ge1\quad\exists C_{H,K}\ge0\quad\exists f_{H,K}:H\to K
\]
hold with \(f_{H,K}\) a \((\lambda_{H,K},C_{H,K})\)-quasi-isometry?

If this holds, can the constants be chosen uniformly over all pairs of horospheres in the fixed universal cover, namely,
\[
 \exists\lambda_X\ge1\quad\exists C_X\ge0\quad\forall H,K\quad\exists f_{H,K}:H\to K
\]
with every \(f_{H,K}\) a \((\lambda_X,C_X)\)-quasi-isometry? Here the constants may depend on \((X,g)\). The source does not explicitly specify the dependence intended by ``uniformly''; this is the natural within-manifold interpretation, and any intended uniformity across different manifolds remains unresolved. The source also gives no explicit range for \(n\).
\subsection{Short English statement}
For a closed negatively curved manifold, must every two horospheres in its universal cover be quasi-isometric when equipped with their intrinsic induced Riemannian distances? If so, can one choose the quasi-isometry constants independently of the two horospheres, for that fixed manifold? The scope of the word ``uniformly'' requires source review.
\subsection{Sources}
\begin{itemize}
\item Local registry: \texttt{data/unsolvedmath/problems.json}, permanent UnsolvedMath ID \(30002138\), source problem code \texttt{catalogue entry}, in the Oberwolfach Reports Open Problems collection. The supplied record lists proposed year 2012 and supplies no individual proposer.
\item Supplied report citation: \emph{Geometric Group Theory, Hyperbolic Dynamics and Symplectic Geometry} (2013), Oberwolfach Reports, pp.~2139--2203, DOI \texttt{10.4171/owr/2012/35}. \url{https://doi.org/10.4171/owr/2012/35}. The citation's year 2013 is retained alongside the supplied proposed year 2012 and the DOI containing 2012; their relationship has not been independently checked.
\item Supplied literature source: Besson--Courtois--Hersonsky, ``horospherical rigidity.'' \url{https://arxiv.org/abs/2411.15093}.
\item Supplied literature source: ``Rigidity of flat holonomies.'' \url{https://doi.org/10.1017/etds.2024.58}.
\item Supplied literature source: ``Oberwolfach Report 35.'' \url{https://doi.org/10.4171/owr/2012/35}.
\item Literature assessment dated \texttt{2026-08-21}: the supplied curation labels the problem open and states that recent horospherical rigidity results impose additional hypotheses and establish special rigidity statements, while quasi-isometry of all horospheres in every closed negatively curved manifold is not established. This is supplied dated curation, not an independently verified literature assessment; time-sensitive status and the cited claims require independent verification.
\item The supplied statement status is \texttt{corrected\_verified}. Its accompanying verification note reports comparison with the official Oberwolfach report and normalization of notation, context, and question form without changing the source-level claim. This verification is reported from the record and has not been independently performed for this draft.
\end{itemize}
\end{document}
